% ID: SEL-D0066-OA-E01
% ESTADO: validada. S066.
\textbf{a)} Multiplicando y sumando entrada a entrada,
\[
A-B+AB^t=\begin{pmatrix}2a+b-1&4-b\\2b&5\end{pmatrix}.
\]
Igualamos a $C$. La entrada inferior izquierda obliga a $b=1$ y entonces $2a=1$. La entrada superior derecha queda en $3$, como se pedía. Así, $\boxed{a=1/2,\ b=1}$.

\textbf{b)} La entrada inferior derecha de $BB^t$ es $0^2+3^2=9$ para todo $b$. Por tanto, el producto nunca es la matriz nula.

\textbf{c)} La igualdad es $AX=-B$. Con $a=1/2$, la segunda fila de $AX$ es dos veces la de $X$, que por tanto vale $(0,-3/2)$. La primera fila de $AX$ es la mitad de la primera fila de $X$ más la segunda. Al igualarla a $(-1,-1)$, resulta
\[
X=\boxed{\begin{pmatrix}-2&1\\0&-3/2\end{pmatrix}}.
\]
